\section{可复现实验路线图}\label{sec:reproduce}

前文的比较依赖两类证据：文献（不可比、有污染）与本文自建回测（样本单一、口径简化）。要把比较推到更牢固的基础上，需要一组\emph{口径统一、可复现、带验证门}的实验。本章给出这样的路线图，按 Phase 0 至 Phase 4 划分，每阶段列明交付物与验证门。

设计原则有三条：\textbf{(i) 变量受控}——两策略共享除信号符号外的全部设定；\textbf{(ii) 口径前置}——成本、样本、指标在实验开始前锁定，不得事后调整；\textbf{(iii) 门禁明确}——每阶段设通过条件，未通过则不得进入下一阶段。

\subsection{Phase 0：基建与口径锁定}

\textbf{交付物}
\begin{itemize}[leftmargin=2.2em]
  \item 统一数据层：同一资产池的日频与分钟频价格、持仓量、成交额；股票需含复权因子与停复牌标记；期货连续合约需含换月跳空守门与前复权处理。
  \item 统一成本层：参与率约束的冲击模型（冲击与参与率平方根成正比，设定参与率上限并对超限标记）；股票端区分佣金、印花税与借券费。
  \item 统一评估层：13 项强制指标（年化收益、年化波动、夏普、最大回撤、Calmar、偏度、超额峰度、年化换手、毛利年化、成本拖累、月胜率、Newey--West $t$ 统计量、成本占毛利比）。
  \item 策略骨架：两套策略共用同一回测引擎，信号函数以符号参数区分方向。
\end{itemize}

\textbf{验证门}
\begin{itemize}[leftmargin=2.2em]
  \item \textbf{G1 一致性}：同一信号、符号翻转后，毛收益序列应严格互为相反数（差异应为零）。这是统一框架的直接检验——若不能通过，说明实现中有非对称处理。
  \item \textbf{G2 成本守恒}：两策略在同一调仓频率下的换手应完全相同；成本拖累应完全相等。若趋势与反转的换手不同，说明信号平滑或阈值处理引入了非对称。
  \item \textbf{G3 无 look-ahead}：在信号计算中注入未来一日数据，回测指标应出现可观测的改善；以此为哨兵检验。
\end{itemize}

G1 与 G2 是本路线图的核心设计。它们在本文自建回测中已经得到验证：期货时序口径下，趋势与反转的换手几乎完全相等（如 $L=5$ 时分别为 52.41 与 52.38），成本拖累一致，净值差异完全来自信号符号——这正是“唯一区别是符号”的实证体现。

\subsection{回测代码框架：五个模块}

Phase 0 的四项交付物落在同一个代码框架里。框架分五层，两套策略共享除信号符号外的全部代码——策略差异被隔离在信号层这一个函数内，G1 与 G2 才能用机器自动校验，而不是靠人工比对。

\begin{quote}\ttfamily\small
数据层：\\
\quad 输入日频行情、复权因子、停复牌标记、期货换月记录\\
\quad 输出统一的收益序列与可交易标记；换月跳空并入收益，停牌日标记为不可交易\\[2pt]
信号层（策略唯一分岔点）：\\
\quad 输入回看窗口 $L$ 与方向符号 $\sigma\in\{+1,-1\}$\\
\quad $s_t \gets \sigma\times\sum_{k=0}^{L-1}\ln(1+r_{t-k})$ \hfill \# 趋势 $\sigma=+1$，反转 $\sigma=-1$\\[2pt]
引擎层（逐日推进，先信号后仓位）：\\
\quad 对每个交易日 $t$：\\
\qquad 用截至 $t$ 日收盘的全部可得信息计算 $s_t$\\
\qquad 目标权重 $w_t \gets$ 缩放与归一化（两策略共用同一函数）\\
\qquad 仓位自 $t+1$ 日起生效，持有期内按收益漂移更新权重\\
\qquad 换手 $\gets\sum_i\bigl|w_{i,t+1}-w^{\mathrm{current}}_{i,t}\bigr|$，仅对换入换出部分计成本\\[2pt]
成本层：\\
\quad 固定基点部分期货单边 2、A 股单边 15；\\
\quad 另设参与率上限 10\% 的冲击项，超限标记该笔不可成交\\[2pt]
评估层：\\
\quad 输出 13 项强制指标；同时输出扣成本前的毛收益序列，供门禁校验\\
\end{quote}

三道哨兵直接挂在引擎层与评估层上，每次跑完自动触发：G1 对 $\sigma=+1$ 与 $\sigma=-1$ 各跑一遍比对毛收益序列，G2 比对两次运行的逐日换手与成本，G3 用信号层的未来数据注入做 look-ahead 探测（定义见 Phase 0）。三道全过，框架才可用于后续实验——任何一道失败，都先修实现再谈结果。

\subsection{Phase 1：尺度扫描与 U 形检验}

\textbf{交付物}：在统一口径下对 $L\in\{5,10,20,60,120,250\}$ 做完整网格，覆盖时序（TS）与截面（XS）两种口径、期货与股票两个市场，输出“$L$—夏普”曲线。

\textbf{验证门}
\begin{itemize}[leftmargin=2.2em]
  \item \textbf{G4 U 形存在性}：至少在一种口径与一个市场上，趋势与反转的相对夏普应随 $L$ 反转（短 $L$ 反转占优、中 $L$ 趋势占优、长 $L$ 反转占优）。若在任何口径上都不出现，则需重新审视“尺度决定归属”的普适性。
  \item \textbf{G5 成本归因}：在同一 $L$ 上，趋势净夏普 $+$ 反转净夏普 应约等于 $-2\times$(成本拖累/波动)。这是成本非对称性的定量检验。
\end{itemize}

\subsection{Phase 2：稳健性与多重检验}

\textbf{交付物}
\begin{itemize}[leftmargin=2.2em]
  \item 样本外检验：滚动 walk-forward，形成期与交易期严格分离。
  \item 多重检验校正：对全部网格组合计算 White Reality Check、Hansen SPA 与紧缩夏普比率；报告校正前后的显著性。
  \item 参数稳健性：$L$ 邻域内的夏普曲面平滑度；是否存在孤立的高夏普点（过拟合征兆）。
\end{itemize}

\textbf{验证门}
\begin{itemize}[leftmargin=2.2em]
  \item \textbf{G6 多重检验后存活}：至少一条“趋势优于反转”或“反转优于趋势”的结论在校正后仍显著。
  \item \textbf{G7 稳健性}：最优 $L$ 的邻域夏普不得出现剧烈跳变；若最优点的夏普显著高于邻域均值，判定为过拟合。
\end{itemize}

\subsection{落地流程：从网格到判决的四道闸门}

把 Phase 2 的第二项交付物展开成可执行的流程。方法学本身（Reality Check、SPA、紧缩夏普）已在第~\ref{sec:audit}~节作为文献标准介绍过，这里只回答操作问题：先算什么、后算什么、哪一步否决哪一步。

\begin{enumerate}[leftmargin=2.2em]
  \item \textbf{锁定试验族}。网格的全部组合（第~\ref{sec:empirical}~节的六窗口 × 两方向 × 四组口径）在实验开始前一次性定义为试验族。族内每个组合都是潜在的“被选中者”，事后追加的组合不得并入——并入即扩大族，先前的校正随之失效。这一步是整个流程的地基，也是最容易被跳过的一步。
  \item \textbf{基准统计量}。对族内每个组合计算净夏普与 Newey--West 调整的 $t$ 统计量（收益自相关下常规标准误失效，见第~\ref{sec:empirical}~节的口径声明），记族内最优者的原始统计量为 $t_{\max}$。到这一步为止，一切与单组合检验无异；$t_{\max}$ 本身不构成证据——它是 48 次搜索的最大值。
  \item \textbf{跨组合校正}。对 $t_{\max}$ 做 block bootstrap 重抽样（保持序列相关的块长），得到 White Reality Check 与 Hansen SPA 的 $p$ 值，回答“最优组合是否只是搜索的产物”；再用 Romano--Wolf 逐步法定位校正后仍显著的组合\emph{集合}——集合可能为空，空集意味着没有任何单一组合在族水平上站得住。
  \item \textbf{稳健性折算与判决}。对通过第 3 步的组合报告紧缩夏普（把估计误差从原始夏普中折掉）与回测过拟合概率（CSCV 口径，0.5 等价于抛硬币）。判决对接 G6：RC 与 SPA 的 $p<0.05$、且过拟合概率明显低于 0.5，判“校正后存活”；任一不满足，该结论降级为方向性描述，只能进比较矩阵的置信度列，不能进结论。
\end{enumerate}

四道闸门的顺序不可调换：先锁定族再算统计量，否则校正对象不确定；先跨组合校正再折算稳健性，否则折算的是一个未经检验的候选。流程跑完输出两张表——校正前后的显著性对照，与被否决组合的清单；后者和前者同样重要，它记录了搜索代价。

\subsection{Phase 3：执行层对照}

\textbf{交付物}
\begin{itemize}[leftmargin=2.2em]
  \item 波动率目标化的不对称检验：对两策略分别施加与不施加波动率目标，比较夏普、偏度与最大回撤的变化。
  \item 资金约束模拟：为反转策略设定保证金缓冲与偏离上限，测试在压力情景（流动性骤降）下的生存性。
  \item 组合融合检验：按尺度分层、按状态加权、按风险预算三种融合方案与简单等权叠加的对比。
\end{itemize}

\textbf{验证门}
\begin{itemize}[leftmargin=2.2em]
  \item \textbf{G8 不对称性确认}：波动率目标化应改善趋势策略的风险调整收益，但不应同等改善（或应恶化）反转策略。若两者同等改善，则第~\ref{sec:execution}~节的核心执行差异判断需修正。
  \item \textbf{G9 融合增益}：至少一种融合方案的夏普或尾部指标优于两策略单独持有与简单叠加。
\end{itemize}

\subsection{Phase 4：沉淀与复现包}

\textbf{交付物}：完整代码仓库（数据层、成本层、引擎、策略、评估），含环境锁定文件、随机种子、以及一键复现脚本；实验记录含每次运行的参数与结果哈希。

\textbf{通过条件}：第三方在无作者协助下，用同一复现包得到同一组指标（允许浮点误差）。

\subsection{与本文实证的关系}

本文第~\ref{sec:empirical}~节的自建回测可以视为本路线图的 Phase 0 与 Phase 1 的初步实现：它通过了 G1、G2（期货时序口径下两策略换手几乎相等），并部分回答了 G4（在期货截面口径上观察到 U 形结构）。Phase 2 至 Phase 4 尚未执行，这是本文实证的未完成部分，也是后续工作的直接入口。
